% Generated draft: requires mathematical and source review.
\documentclass[11pt]{article}
\usepackage[margin=25mm]{geometry}
\usepackage{fontspec}
\setmainfont{Latin Modern Roman}
\usepackage{amsmath,amssymb,mathtools,unicode-math}
\setmathfont{Latin Modern Math}
\usepackage{enumitem,needspace,xurl,hyperref}
\setlength{\parindent}{0pt}
\setlength{\parskip}{5pt}
\setlength{\emergencystretch}{3em}
\newcommand{\N}{\mathbb{N}}
\newcommand{\Z}{\mathbb{Z}}
\newcommand{\Q}{\mathbb{Q}}
\newcommand{\R}{\mathbb{R}}
\newcommand{\C}{\mathbb{C}}
\newcommand{\F}{\mathbb{F}}
\newcommand{\A}{\mathbb{A}}
\newcommand{\PP}{\mathbb{P}}
\newcommand{\E}{\mathbb{E}}
\newcommand{\eps}{\varepsilon}
\newcommand{\dd}{\,\mathrm{d}}
\begin{document}
\section*{UnsolvedMath ID: 20003100}
\textbf{Source problem code:} catalogue entry
\subsection{Definitions and mathematical statement}
The source asks whether a stretch line is typically recurrent, dense, or equidistributed in moduli space. Its word ``typically'' does not choose a measure or parameterization. The supplied result specifies closed oriented surfaces of genus $g\geq2$, directed classical unit-speed Thurston stretch lines, and normalized Mirzakhani measure on the unit-length measured-lamination phase space
\[
\mathcal P^1\mathcal M_g=\{(X,\lambda):X\in\mathcal T_g,\ \lambda\in\mathcal{ML}(S)\setminus\{0\},\ \ell_X(\lambda)=1\}/\operatorname{Mod}(S).
\]
Let $\mu_{\mathrm{Th}}$ be Thurston measure on measured laminations and define
\[
B(X)=\mu_{\mathrm{Th}}\{\lambda:\ell_X(\lambda)\leq1\},\qquad
b_g=\int_{\mathcal M_g}B(X)\,\dd\mu_{\mathrm{WP}}(X).
\]
The phase probability $\widehat\mu_{\mathrm{Mir}}$ uses the radial pushforward $\lambda\mapsto\lambda/\ell_X(\lambda)$ of the length-at-most-one cone in each fiber, together with Weil--Petersson measure, normalized by $b_g$. Thus its moduli-space marginal is
\[
\dd\nu_g(X)=\frac{B(X)}{b_g}\,\dd\mu_{\mathrm{WP}}(X).
\]
In particular the asserted equilibrium law is not simply normalized Weil--Petersson measure.

The aid invokes Calderon--Farre's ergodicity and geodesicity results to identify a full-measure domain of a measure-preserving ergodic phase flow $\Phi_t$ whose projected orbits $X_t$ are these classical directed stretch lines. With exactly those theorem hypotheses and the unit-speed time parameter, it concludes that for $\widehat\mu_{\mathrm{Mir}}$-almost every initial phase point the phase orbit is recurrent, its moduli projection is dense, and
\[
\frac1T\int_0^T\delta_{X_t}\,\dd t\ \Longrightarrow\ \nu_g\qquad(T\to\infty).
\]
Equivalently the averages of bounded continuous test functions converge to their $\nu_g$ integrals. The common almost-everywhere statement uses the specified phase measure; it is not a claim for every starting surface, every fixed lamination, or every stretch line.

The aid's stronger visit conclusion is that, simultaneously for every nonempty moduli-space open set $U$,
\[
\liminf_{T\to\infty}\frac1T\operatorname{Leb}\{t\in[0,T]:X_t\in U\}>0.
\]
It follows from ergodicity on a countable base and full support of $\nu_g$, without requiring zero boundary measure for arbitrary $U$. Consequently $\liminf_{t\to\infty}\operatorname{sys}(X_t)=0$, while for some fixed $\epsilon_0>0$ the set of times with $\operatorname{sys}(X_t)\geq\epsilon_0$ has positive lower frequency. Arbitrarily thin excursions and repeated thick-part returns are compatible. Recurrence is a phase-flow conclusion from Poincar\'e recurrence; equidistribution follows from Birkhoff and the stated marginal formula. The precise phase-flow construction, full-measure domain, and theorem references are not included in the supplied record.

The trailing source sentence about Lenzhen--Masur concerns when Teichm\"uller geodesic rays remain bounded distance apart. It is retained as separate context and is not a recurrence theorem for Thurston stretch lines.
\subsection{Short English statement}
With a specified closed-surface phase measure and unit-speed stretch flow, the supplied aid reports almost-everywhere recurrence, density, and equidistribution. The projected probability has density $B/b_g$ relative to Weil--Petersson measure. Typical lines visit every open set with positive lower frequency, including thin regions and a fixed thick part. This is a measure-specific statement, not one about all stretch lines.
\subsection{Sources}
\begin{itemize}
\item Local registry: permanent catalogue entry, exact source code catalogue entry; \texttt{aim-topology-notes.json}, \texttt{notes[187]}.
\item Weixu Su's 2014 AIM list \emph{Problems on Thurston Metric}, Problem 2.10 attributed to K. Rafi; workshop catalog title \emph{Lipschitz metric on Teichmueller space}, \url{https://aimath.org/pastworkshops/lipschitzteichproblems.pdf}.
\item Supplied attributions: Calderon--Farre for stretch-flow ergodicity and geodesicity; Mirzakhani for the projection formula. Poincar\'e recurrence and Birkhoff's theorem are named inputs. No full paper titles, theorem numbers, or complete authoritative citations are included.
\item The trailing source names Lenzhen and Masur for a separate bounded-distance question about Teichm\"uller geodesic rays, without a full citation.
\item Supplied research attempt 1, registry updated 14 August 2026: partial result, low novelty confidence, machine-generated and not peer reviewed. Statement status is \texttt{reconstructed\_unverified}, clean statement is null, and verification reports inherited PDF inspection but ends before reproducing a fully checked question.
\end{itemize}
\end{document}
